% ============================================================
% EEEC 1099 Biomedical Sensors - Spring 2026
% Final Project Proposal
% Due: 2026-05-22 (email to EdwinKan5@nycu.edu.tw)
% ============================================================
\documentclass[11pt,a4paper]{article}

\usepackage[margin=0.9in]{geometry}
\usepackage{parskip}
\usepackage{cite}
\usepackage[hidelinks]{hyperref}

\setlength{\parskip}{0.35em}
\setlength{\parindent}{0pt}

% Custom enumerate spacing without enumitem dependency
\newenvironment{tightenum}{%
  \begin{enumerate}%
    \setlength{\itemsep}{0.4em}%
    \setlength{\parsep}{0pt}%
    \setlength{\topsep}{0.2em}%
    \setlength{\leftmargin}{1.5em}%
}{\end{enumerate}}

\begin{document}

% --------- Title block ---------
\begin{center}
{\Large\bfseries When Should the LF/HF Ratio Be Used as an\\
Autonomic-Balance Metric in Wearable ECG?\par}
\vspace{0.5em}
Pei-En Wu\\
{\small NYCU ID: 113511103 \quad
Email: \href{mailto:peienwu.ee13@nycu.edu.tw}{peienwu.ee13@nycu.edu.tw}\par}
{\small EEEC 1099 Introduction to Biomedical Sensors, Spring 2026 --- Final Project Proposal}
\end{center}

\vspace{0.5em}

% --------- Three major points ---------
\section*{Three Major Points}

\begin{tightenum}

\item This project focuses on LF/HF, a frequency-domain HRV ratio that
emerged from classic spectral-analysis work as a simplified marker of
sympathovagal balance~\cite{taskforce1996},~\cite{malliani1991}, but later
cardiovascular physiology literature questioned whether LF and HF power can
cleanly separate sympathetic and parasympathetic
activity~\cite{eckberg1997},~\cite{billman2013}.

\item Using fourteen single-lead ECG recordings from my HW5 dataset, I will
use posture, breath-hold, walking, and paced-breathing conditions as an
illustrative within-subject case study to discuss how respiration, posture,
motion, and window duration can reshape LF/HF in short-term wearable
recordings~\cite{bernardi2000}, especially when the respiratory peak crosses
the 0.15-Hz LF/HF boundary during slow paced breathing~\cite{billman2013}.

\item By synthesizing HRV standards~\cite{taskforce1996}, the original LF/HF
framework~\cite{malliani1991}, methodological
critiques~\cite{eckberg1997,bernardi2000,billman2013}, and reporting
recommendations~\cite{laborde2017}, the final deliverable will be a
context-aware HRV reporting checklist that, for each common short-term wearable
recording context, specifies the primary metric, its interpretive validity, and
the dominant caveat.

\end{tightenum}

% --------- References ---------
% thebibliography emits its own ``References'' heading (\refname); do not add \section*{References}.

\begin{thebibliography}{6}
\setlength{\itemsep}{0.15em}

\bibitem{taskforce1996}
Task Force of the European Society of Cardiology and the North American
Society of Pacing and Electrophysiology,
``Heart rate variability: Standards of measurement, physiological
interpretation, and clinical use,''
\emph{Circulation}, vol.~93, no.~5, pp.~1043--1065, 1996,
doi: \href{https://doi.org/10.1161/01.CIR.93.5.1043}{10.1161/01.CIR.93.5.1043}.

\bibitem{malliani1991}
A.~Malliani, M.~Pagani, F.~Lombardi, and S.~Cerutti,
``Cardiovascular neural regulation explored in the frequency domain,''
\emph{Circulation}, vol.~84, no.~2, pp.~482--492, 1991,
doi: \href{https://doi.org/10.1161/01.CIR.84.2.482}{10.1161/01.CIR.84.2.482}.

\bibitem{eckberg1997}
D.~L. Eckberg,
``Sympathovagal balance: A critical appraisal,''
\emph{Circulation}, vol.~96, no.~9, pp.~3224--3232, 1997,
doi: \href{https://doi.org/10.1161/01.CIR.96.9.3224}{10.1161/01.CIR.96.9.3224}.

\bibitem{bernardi2000}
L.~Bernardi et al.,
``Effects of controlled breathing, mental activity and mental stress with or
without verbalization on heart rate variability,''
\emph{Journal of the American College of Cardiology}, vol.~35, no.~6,
pp.~1462--1469, 2000,
doi: \href{https://doi.org/10.1016/S0735-1097(00)00595-7}{10.1016/S0735-1097(00)00595-7}.

\bibitem{billman2013}
G.~E. Billman,
``The LF/HF ratio does not accurately measure cardiac sympatho-vagal
balance,''
\emph{Frontiers in Physiology}, vol.~4, art.~26, 2013,
doi: \href{https://doi.org/10.3389/fphys.2013.00026}{10.3389/fphys.2013.00026}.

\bibitem{laborde2017}
S.~Laborde, E.~Mosley, and J.~F. Thayer,
``Heart rate variability and cardiac vagal tone in psychophysiological
research: Recommendations for experiment planning, data analysis, and data
reporting,''
\emph{Frontiers in Psychology}, vol.~8, art.~213, 2017,
doi: \href{https://doi.org/10.3389/fpsyg.2017.00213}{10.3389/fpsyg.2017.00213}.

\end{thebibliography}

\end{document}
